\section*{The Propers: Themes and Movement}

\begin{massbox}{A people received, reconciled, and sent}
The Lectionary joins Israel's covenant vocation to Christ's compassionate commissioning of the Twelve, while Romans supplies an independent apostolic proclamation of reconciliation. The Missal frames this hearing with grace for action, bodily and spiritual renewal, and sacramental unity; the movement is a synthesis of texts the liturgy places together: the Lectionary pairs Exodus with Matthew for Year A, Romans keeps its own course, and the Week XI orations serve every Sunday cycle.
\end{massbox}

\begin{center}
\textbf{Called by grace $\longrightarrow$ gathered as God's people $\longrightarrow$ sent from Christ's compassion $\longrightarrow$ renewed in Communion}
\end{center}

\subsection*{1. Need becomes petition\enspace\elementref{Entrance, Collect}}

The Entrance does not begin with institutional confidence but with a plea to be heard and helped. The Collect likewise acknowledges mortal frailty and asks grace for both willing and doing. Mission therefore begins as dependence: the Church cannot make herself the Lord's harvest, manufacture vocation, or turn resolve into holy action apart from God.

\subsection*{2. Election becomes vocation\enspace\elementref{First Reading, Psalm, Gospel}}

Exodus first announces what God has done in bearing Israel from bondage; only then does it name covenant obedience and priestly vocation. Psalm 100 answers with belonging and grateful service. Matthew's official correlation remains within Israel's covenant history: Jesus sends the Twelve first to Israel, and the mission remains governed by the Messiah of Israel and Lord of the harvest.

\subsection*{3. Compassion becomes gratuitous mission\enspace\elementref{Acclamation, Gospel, Second Reading}}

Jesus' compassion sees crowds as harassed and abandoned, then commands prayer and gives authority. Chrysostom stresses that the harvest was already prepared and remains Christ's. “Without cost” disciplines pride and covetousness. Romans independently supplies the deepest prior gift: Christ acts for the weak, sinful, and hostile, so Christian service gives from reconciliation already received.

\subsection*{4. Offering becomes renewal and unity\enspace\elementref{Prayer over Offerings, Communion, Prayer after Communion}}

The Prayer over the Offerings holds body and mind together: ordinary food sustains creaturely life and Sacrament renews it. Communion may voice desire for God's house or Christ's prayer for unity. The concluding prayer names unity as both sign and effect. These Week XI texts can frame Year A mission, though every Sunday cycle shares them.

\begin{studybox}{Israel, ministry, and Romans}
\textbf{Israel:} Covenant vocation and the initial mission remain within Israel's salvation history. \textbf{Ministry:} gratuitous grace does not deny lawful material support, but it forbids treating divine gifts as personal property. \textbf{Relationship:} Romans remains an independent semi-continuous strand.
\end{studybox}
